\section{Edge-half for sufficiently large graphs at half density}
\label{sec:four-edge-half-dense-small-cuts}

\begin{lemma}\label{lem:four-full-pole-short-cycle}
Let $X$ be a connected normalized $k$-pole of rank $h$, with $k\ge2$.
If $h>5k$ and some terminal pair has $2^h$ paths, then $X$ contains a cycle of length at most four.
\end{lemma}

\begin{proof}
Suppress the $k-2$ unused ports and use the saturated two-pole classification.
Its positive-rank pieces are digons or capped singleton-and-square ladders.
The candidate construction in Theorem~\ref{thm:saturated-pair-small-cyclic-cut} gives at least $h/5$ vertex-disjoint candidates.
For a piece of rank at most four, its whole vertex set is a candidate.
It contains a digon, a cap triangle when its first block is a singleton, or a cap four-cycle when its first block is a square.
For a larger piece, each candidate is a square block or two consecutive singleton blocks, and contains a four-cycle.

Restoring the unused ports subdivides at most $k-2$ edges.
Each restored port can spoil at most one candidate, since their vertex sets are disjoint.
An unspoiled candidate therefore retains a cycle of length at most four in $X$.
Such a candidate exists because $h/5>k>k-2$.
\end{proof}

\begin{lemma}\label{lem:four-girth-five-small-pole-ranks}
A simple normalized $k$-pole of girth at least five has rank at least four when $k=3$, and at least three when $k=4$.
\end{lemma}

\begin{proof}
Let $n$ denote its order.
There is at most one path of length one or two between each ordered pair of distinct vertices.
The numbers of such paths sum to
\[
 (3n-k)+(6n-4k)=9n-5k.
\]
Thus $n(n-1)\ge9n-5k$, or $n^2-10n+5k\ge0$.
Use $n=2h+k-2$ and $h\ge1$.
For $k=3$, this forces $n\ge9$ and $h\ge4$.
For $k=4$, it forces $n\ge8$ and $h\ge3$.
\end{proof}

\begin{proposition}\label{prop:four-girth-five-small-cut-deficit}
For all sufficiently large ranks, every three-edge-connected simple cubic graph of girth at least five with a cyclic cut of size at most four satisfies
\[
 c(G)<2^{\beta(G)-1}.
\]
\end{proposition}

\begin{proof}
Choose a cyclic cut of minimum size $j\in\{3,4\}$.
Its shores are connected normalized $j$-poles with distinct ports.
Indeed, disconnected shores or repeated attachments would produce a smaller cyclic cut, or an edge cut of size at most two.
Write the shores as $X,Y$, with ranks $h\le k$.
Then $\beta(G)=h+k+j-1$, and $k$ tends to infinity with $\beta(G)$.
Lemma~\ref{lem:four-full-pole-short-cycle} excludes every saturated terminal pair in $Y$ for sufficiently large $k$.
The path saturation gap therefore bounds each normalized path coordinate of $Y$ by $15/16$.

Suppose first that $j=3$.
If two path coordinates exceed $0.51$, the two-high-pair theorem makes the third $O(2^{-k/2})$.
Their sum is then at most $15/8+O(2^{-k/2})$.
Otherwise, their sum is at most
\[
 \frac{15}{16}+2(0.51)=\frac{783}{400}=1.9575.
\]
Each path coordinate of $X$ is at most one.
The exact three-cut identity, the trivial bound $c(X)<2^h$, and the asymptotic cycle bound give
\[
 \frac{c(G)}{2^{h+k}}
 \le \frac{783}{400}+o(1)+\frac9{16}2^{-h}+2^{-k}
 \le \frac{783}{400}+\frac9{256}+o(1)<2.
\]
Here $h\ge4$ by Lemma~\ref{lem:four-girth-five-small-pole-ranks}.
Every error is uniform in $h\le k$.

Now suppose that $j=4$.
Use the normalized state notation $S(Y),F(Y)$ from the four-cut rigidity proof.
If all six path coordinates of $Y$ are at most $0.51$, the low-state inequality gives $S(Y)+F(Y)\le3.82$.
Since $h\ge3$, the exact gluing formula yields
\[
 \frac{c(G)}{2^{h+k}}
 \le3.82+\frac9{16}2^{-h}+2^{-k}
 \le3.82+\frac9{128}+2^{-k}<4.
\]
If a path coordinate exceeds $0.51$, the high-path cycle tradeoff gives $c(Y)=O(2^{k/2})$.
The aggregate rigidity lemma forces $S(Y)+F(Y)<3.99$, since $Y$ has no saturated pair.
Consequently
\[
 \frac{c(G)}{2^{h+k}}
 \le3.99+O(2^{-h-k/2})+2^{-k}<4
\]
uniformly for $h\ge3$ and sufficiently large $k$.
This proves both cases.
\end{proof}

\begin{corollary}\label{cor:four-dense-edge-half}
There is an absolute rank threshold $R$ such that every connected simple cubic graph $G$ of rank $r\ge R$ with
\[
 c(G)\ge2^{r-1}
\]
has an edge $e$ satisfying $2c_e(G)\le c(G)$.
\end{corollary}

\begin{proof}
The existing local identities handle bridges, two-edge cuts, triangles, and four-cycles.
Thus assume that $G$ is three-edge-connected and has girth at least five.
Proposition~\ref{prop:four-girth-five-small-cut-deficit} excludes a cyclic cut of size at most four at sufficiently large rank.
The established below-half reduction without small cyclic cuts, Corollary~\ref{cor:new-edge-half-below-density}, now supplies the desired edge.
\end{proof}

This result does not settle edge-half below half density or at the finitely many ranks below the threshold.
